% ----------------------------------------------------------
% Pacotes básicos 
% ----------------------------------------------------------
% \usepackage{Pacotes/abntex2}
\usepackage{amsmath, amssymb, amsfonts, amsthm, mathtools}
\usepackage{mathptmx} 							% Usa a fonte Times New Roman	 (UDESC/CCT)
\usepackage[T1]{fontenc}						% Seleção de códigos de fonte.
\usepackage[utf8]{inputenc}						% Codificação do documento (conversão automática dos acentos)
\usepackage{lastpage}							% Usado pela Ficha catalográfica
\usepackage{indentfirst}						% Indenta o primeiro parágrafo de cada seção.
\usepackage[dvipsnames,table]{xcolor}			% Controle das cores
\usepackage{graphicx}							% Inclusão de gráficos
\usepackage{microtype} 							% para melhorias de justificação
% \usepackage{lipsum}								% para geração de dummy text
\usepackage[brazilian,hyperpageref]{backref}	% Paginas com as citações na bibl
\usepackage[alf,abnt-emphasize=bf,abnt-full-initials=yes]{abntex2cite}					% Citações padrão ABNT
\usepackage{adjustbox}							% Pacote de ajuste de boxes
\usepackage{subcaption}							% Inclusão de Subfiguras e sublegendas		
\usepackage{enumitem}							% Personalização de listas
\usepackage[section]{placeins}					% Manter as figuras delimitadas na respectiva seção com a opção [section]
\usepackage{multirow}							% Multi colunas nas tabelas
\usepackage{array,tabularx} 					% Pacotes de tabelas
\usepackage{booktabs}							% Pacote de tabela profissional
\usepackage{rotating}							% Rotacionar figuras e tabelas
\usepackage{xfrac}								% Fazer frações n/d em linha
\usepackage{bm}									% Negrito em modo matemático
\usepackage{xstring}							% Manipulação de strings
\usepackage{chngcntr}							% Pacte usado para deixar numeração de equações sequencial (UDESC/CCT)
\usepackage{xspace}								% Espaçamento após comandos
\usepackage{bussproofs}							% Provas em estilo de árvore
\usepackage{hyperref}							% Suporte para hiperlinks
\usepackage{listings}					% Highlighting de código Coq
% \usepackage{llvm/lang}  % include custom language for LLVM IR.
% \usepackage{nasm/style} % include custom style for NASM assembly.
% \usepackage{float}	
\lstset{
  basicstyle=\ttfamily,
  mathescape
}
\usepackage{fancyvrb}
\usepackage{stmaryrd}
\usepackage{minted}

\theoremstyle{definition}
\newtheorem{definition}{Definição}[section]

\theoremstyle{remark}
\newtheorem*{remark}{Remark}

\definecolor{darkgreen}{rgb}{0,0.6,0}
\definecolor{darkorange}{rgb}{0.8,0.4,0.2}
\definecolor{darkred}{rgb}{0.8,0,0}

\counterwithout{equation}{chapter}
% fonte: https://latex.org/forum/viewtopic.php?t=15392

% Comando para deixar numeração das equações contínua (1), (2), (3)... ao invés de organizar por capítulos (1.1)(1.2)... (2.1)(2.2)
%\renewcommand{\theequation}{\arabic{equation}}

%\numberwithin{equation}{section}

% Cabeçalho somente com numeração de página 10pt
\makepagestyle{PagNumReduzida}
\makeevenhead{PagNumReduzida}{\ABNTEXfontereduzida\thepage}{}{}
\makeoddhead{PagNumReduzida}{}{}{\ABNTEXfontereduzida\thepage}
%fonte: https://github.com/abntex/abntex2/wiki/HowToCustomizarCabecalhoRodape
%fonte: Manual memoir seção 7.3 pg. 111 pdf http://linorg.usp.br/CTAN/macros/latex/contrib/memoir/memman.pdf 

% Personalização das opções das listas
\setlist[itemize]{leftmargin=\parindent}

% Citação online --- MODIFICAR ---
\newcommand{\citeshort}[1]{\citeauthoronline{#1}~(\citeyear{#1})}

\newcommand{\me}[1]{Elaborado pelo autor (#1).}

% Configuração do pgfplots
% \pgfplotsset{compat=newest} %compat=1.14
% \pgfplotsset{plot coordinates/math parser=false} 
% \newlength\figureheight 
% \newlength\figurewidth 

% Libraries do TiKz
% \usetikzlibrary{quotes,angles,arrows}
% \usetikzlibrary{through,calc,math}
% \usetikzlibrary{graphs,backgrounds,fit}
% \usetikzlibrary{shapes,positioning,patterns,shadows}
% \usetikzlibrary{decorations.pathreplacing}
% \usetikzlibrary{shapes.geometric}
% \usetikzlibrary{arrows.meta}
% \usetikzlibrary{external}

%\tikzexternalize[]
%\tikzexternalenable
%\tikzexternalize
%\tikzexternaldisable
%\tikzset{external/force remake}
%\tikzexternalize[shell escape=-enable-write18]

% Personalização das legendas
\usepackage[format = plain, %hang
			justification = centering,
			labelsep = endash,
			singlelinecheck = false,
			skip = 6pt,
			listformat = simple]{caption}	

% Personalizações de tipo de colunas de tabelas
\newcolumntype{L}[1]{>{\raggedright\let\newline\\\arraybackslash\hspace{0pt}}m{#1}}
\newcolumntype{C}[1]{>{\centering\let\newline\\\arraybackslash\hspace{0pt}}m{#1}}
\newcolumntype{R}[1]{>{\raggedleft\let\newline\\\arraybackslash\hspace{0pt}}m{#1}}

% Personalizações de cores da UDESC
\definecolor{CapaAmareloUDESC}{RGB}{243,186,83}		% Especialização
\definecolor{CapaVerdeUDESC}{RGB}{0,112,52}			% Mestrado
\definecolor{CapaVermelhoUDESC}{RGB}{171,35,21}		% Doutorado
\definecolor{CapaAzulUDESC}{RGB}{38,54,118} 		% Pós-Doutorado

% CONFIGURAÇÕES DE PACOTES
% Configurações do pacote backref
% Usado sem a opção hyperpageref de backref
\renewcommand{\backrefpagesname}{Citado na(s) página(s):~}
% Texto padrão antes do número das páginas
\renewcommand{\backref}{}
% Define os textos da citação
\renewcommand*{\backrefalt}[4]{
	\ifcase #1 %
	Nenhuma citação no texto.%
	\or
	Citado na página #2.%
	\else
	Citado #1 vezes nas páginas #2.%
	\fi}%

% alterando o aspecto da cor azul
%\definecolor{blue}{RGB}{41,5,195}

% informações do PDF
\makeatletter
\hypersetup{
	%pagebackref=true,
	pdftitle={\@title}, 
	pdfauthor={\@author},
	pdfsubject={\imprimirpreambulo},
	pdfcreator={LaTeX with abnTeX2},
	pdfkeywords={abnt}{latex}{abntex}{abntex2}{trabalho academico}, 
	colorlinks=true,       		% false: boxed links; true: colored links
	linkcolor=black,          	% color of internal links
	citecolor=black,        	% color of links to bibliography
	filecolor=black,      		% color of file links
	urlcolor=black,
	bookmarksdepth=4
}
\makeatother


% \makeatletter
% \newcommand{\includetikz}[1]{%
% 	\tikzsetnextfilename{#1}%
% 	\input{#1.tex}%
% }
% \makeatother

% ---
% Possibilita criação de Quadros e Lista de quadros.
% Ver https://github.com/abntex/abntex2/issues/176
%
\newcommand{\quadroname}{Quadro}
\newcommand{\listofquadrosname}{Lista de quadros}

\newlistof{listofquadros}{loq}{\listofquadrosname}
\newlistentry{quadro}{loq}{0}

% configurações para atender às regras da ABNT
\setfloatadjustment{quadro}{\centering}
\counterwithout{quadro}{chapter}
\renewcommand{\cftquadroname}{\quadroname\space} 
\renewcommand*{\cftquadroaftersnum}{\hfill--\hfill}

\setfloatlocations{quadro}{hbtp} % Ver https://github.com/abntex/abntex2/issues/176
% ---


% Espaçamento depois do título
\setlength{\afterchapskip}{0.7\baselineskip}
% O tamanho do parágrafo é dado por:
\setlength{\parindent}{1.25cm}
% Controle do espaçamento entre um parágrafo e outro:
\setlength{\parskip}{0.0cm}  % tente também \onelineskip
%\SingleSpacing % Espaçamento simples 
\OnehalfSpacing % Espaçamento 1,5 (UDESC/CCT)
%\DoubleSpacing	% Espaçamento duplo

% ---
% Margens - NBR 14724/2011 - 5.1 Formato
% ---
\setlrmarginsandblock{3cm}{2cm}{*}
\setulmarginsandblock{3cm}{2cm}{*}
\checkandfixthelayout[fixed]
% ---


% To use externalize consider
%https://tex.stackexchange.com/questions/182783/tikzexternalize-not-compatible-with-miktex-2-9-abntex2-package
%Lauro Cesar digged into the problem until he came with a solution for me to test. And it Works!
%
%According to this link:
%
%The package calc changed the commands \setcounter and friends to be fragile. 
% So you have to make them robust. The example below uses etoolbox with \robustify:
%
\usepackage{etoolbox}
\robustify\setcounter
\robustify\addtocounter
\robustify\setlength
\robustify\addtolength


%% How to silence memoir class warning against the use of caption package?
%% https://tex.stackexchange.com/questions/391993/how-to-silence-memoir-class-warning-against-the-use-of-caption-package
%\usepackage{silence}
%\WarningFilter*{memoir}{You are using the caption package with the memoir class}
%\WarningFilter*{Class memoir Warning}{You are using the caption package with the memoir class}

% --------------------------------------------------------
% INICIO DAS CUSTOMIZACOES PARA A UDESC
% --------------------------------------------------------

% --------------------------------------------------------
% Fontes padroes de part, chapter, section, subsection e subsubsection
% --------------------------------------------------------
% --- Chapter ---
\renewcommand{\ABNTEXchapterfont}{\fontseries{b}} %\bfseries
\renewcommand{\ABNTEXchapterfontsize}{\normalsize}
% --- Part ---
\renewcommand{\ABNTEXpartfont}{\ABNTEXchapterfont}
\renewcommand{\ABNTEXpartfontsize}{\LARGE}
% --- Section ---
\renewcommand{\ABNTEXsectionfont}{\normalfont}
\renewcommand{\ABNTEXsectionfontsize}{\normalsize}
% --- SubSection ---
\renewcommand{\ABNTEXsubsectionfont}{\fontseries{b}} %\bfseries
\renewcommand{\ABNTEXsubsectionfontsize}{\normalsize}
% --- SubSubSection ---
\renewcommand{\ABNTEXsubsubsectionfont}{\itshape}
\renewcommand{\ABNTEXsubsubsectionfontsize}{\normalsize}

\renewcommand{\ABNTEXsubsubsubsectionfont}{\normalfont}
\renewcommand{\ABNTEXsubsubsubsectionfontsize}{\normalsize}
% ---

% --------------------------------------------------------
% Fontes das entradas do sumario
% --------------------------------------------------------

\renewcommand{\cftpartfont}{\ABNTEXpartfont\selectfont}
\renewcommand{\cftpartpagefont}{\normalsize\selectfont}

\renewcommand{\cftchapterfont}{\ABNTEXchapterfont\selectfont}
\renewcommand{\cftchapterpagefont}{\normalsize\selectfont}

\renewcommand{\cftsectionfont}{\ABNTEXsectionfont\selectfont}
\renewcommand{\cftsectionpagefont}{\normalsize\selectfont}

\renewcommand{\cftsubsectionfont}{\ABNTEXsubsectionfont\selectfont}
\renewcommand{\cftsubsectionpagefont}{\normalsize\selectfont}

\renewcommand{\cftsubsubsectionfont}{\normalfont\itshape\selectfont}
\renewcommand{\cftsubsubsectionpagefont}{\normalsize\selectfont}

\renewcommand{\cftparagraphfont}{\normalfont\selectfont}
\renewcommand{\cftparagraphpagefont}{\normalsize\selectfont}

% --------------------------------------------------------
% Usando os pacotes hyperref, uppercase... 
% Para deixar a section do toc uppercase precisa de:
% --------------------------------------------------------
\usepackage{textcase}

\makeatletter

\let\oldcontentsline\contentsline
\def\contentsline#1#2{%
	\expandafter\ifx\csname l@#1\endcsname\l@section
	\expandafter\@firstoftwo
	\else
	\expandafter\@secondoftwo
	\fi
	{%
		\oldcontentsline{#1}{\MakeTextUppercase{#2}}%
	}{%
		\oldcontentsline{#1}{#2}%
	}%
}
\makeatother

% --------------------------------------------------------
% Renomenando as entradas de APÊNDICES E ANEXOS
% --------------------------------------------------------

\renewcommand{\apendicesname}{AP\^ENDICES}
\renewcommand{\anexosname}{ANEXOS}


% Manipulação de Strings
%\RequirePackage{xstring}

% Comando para inverter sobrenome e nome
\newcommand{\invertname}[1]{%
	\StrBehind{#1}{{}}, \StrBefore{#1}{{}}%
}%


% --------------------------------------------------------
% Alterando os estilos de Caption e Fonte
% --------------------------------------------------------
\makeatletter
% Define o comando \fonte que respeita as configurações de caption do memoir ou do caption
\renewcommand{\fonte}[2][\fontename]{%
	\M@gettitle{#2}%
	\memlegendinfo{#2}%
	\par
	\begingroup
	\@parboxrestore
	\if@minipage
	\@setminipage
	\fi
	\ABNTEXfontereduzida
	\configureseparator
	\captiondelim{\ABNTEXcaptionfontedelim}
	\@makecaption{#1}{\ignorespaces #2}\par
	\endgroup}


\captionstyle[\raggedright]{\raggedright}

\makeatother

\setlength{\cftbeforechapterskip}{0pt plus 0pt}
\renewcommand*{\insertchapterspace}{}

\newlength{\mylen}	% New length to use with spacing
\setlength{\mylen}{1pt}

\setlength{\cftbeforechapterskip}{\mylen}
\setlength{\cftbeforesectionskip}{\mylen}
\setlength{\cftbeforesubsectionskip}{\mylen}
\setlength{\cftbeforesubsubsectionskip}{\mylen}
\setlength{\cftbeforesubsubsubsectionskip}{\mylen}


% ---
% Ajuste das listas de abreviaturas e siglas ; e símbolos [Personalizada para UDESC com espaçamento 1,5]
% ---

% ---
% Redefinição da Lista de abreviaturas e siglas [Personalizada para UDESC com espaçamento 1,5]
\renewenvironment{siglas}{%
	\pretextualchapter{\listadesiglasname}
	\begin{symbols} 
		\setlength{\itemsep}{0pt}	% Ajuste para Espaçamento 1,5 (UDESC/CCT)
	}{% 
	\end{symbols}
	\cleardoublepage
}
% ---

% ---
% Redefinição da Lista de símbolos [Personalizada para UDESC com espaçamento 1,5]
\renewenvironment{simbolos}{%
	\pretextualchapter{\listadesimbolosname}
	\begin{symbols}
		\setlength{\itemsep}{0pt}	% Ajuste para Espaçamento 1,5 (UDESC/CCT)
	}{%
	\end{symbols}
	\cleardoublepage
}
% ---

\usepackage{xcolor}
\usepackage{textcomp}

% TODO: Add support for digits.
%\definecolor{digits}{RGB}{0,0,128} % dark blue

\definecolor{comment}{RGB}{0,128,0}     % dark green
\definecolor{string}{RGB}{255,0,0}      % red
\definecolor{instruction}{RGB}{0,0,255} % blue
\definecolor{directive}{RGB}{128,0,128} % purple
\definecolor{register}{RGB}{128,0,0}    % dark red
\definecolor{darker}{RGB}{0, 0, 0}

\lstdefinestyle{nasm}{
	commentstyle=\color{comment},
	stringstyle=\color{string},
	keywordstyle=\color{instruction},
	keywordstyle=[2]\color{directive},
	keywordstyle=[3]\color{register},
	basicstyle=\footnotesize\ttfamily,
	% numbers=left,
	numberstyle=\tiny,
	numbersep=5pt,
	frame=none,
	breaklines=true,
	prebreak=\raisebox{0ex}[0ex][0ex]{\ensuremath{\hookleftarrow}},
	showstringspaces=false,
	upquote=true,
	tabsize=8,
}

\lstdefinestyle{asdf}{
	keywordstyle=\bfseries,
	basicstyle=\footnotesize\ttfamily,
	% numbers=left,
	numberstyle=\tiny,
	numbersep=5pt,
	frame=none,
	breaklines=true,
	prebreak=\raisebox{0ex}[0ex][0ex]{\ensuremath{\hookleftarrow}},
	showstringspaces=false,
	upquote=true,
	tabsize=8,
}

\lstdefinelanguage{mygramar}{
        sensitive=true,
	alsoletter={\%},
	keywords=[1]{def},
}

\lstdefinelanguage{llvm}{
	sensitive=true,
	alsoletter={\%},
	% comments.
	%    ; line comment
	comment=[l]{;},
	% strings.
	%    "foo"
	string=[b]{"},
	% instructions.
	%    ref: http://llvm.org/docs/LangRef.html#instruction-reference
	keywords=[1]{
add, addrspacecast, alloca, and, ashr, atomicrmw, bitcast, br, call, cmpxchg,
extractelement, extractvalue, fadd, fcmp, fdiv, fence, fmul, fpext, fptosi,
fptoui, fptrunc, frem, fsub, getelementptr, icmp, indirectbr, insertelement,
insertvalue, inttoptr, invoke, landingpad, load, lshr, mul, or, phi, ptrtoint,
resume, ret, sdiv, select, sext, shl, shufflevector, sitofp, srem, store, sub,
switch, to, trunc, udiv, uitofp, unreachable, urem, va_arg, xor, zext
	},
	% directives.
	%    ref: http://llvm.org/docs/LangRef.html
	keywords=[2]{
acq_rel, acquire, addrspace, alias, align, alignstack, alwaysinline, any,
anyregcc, appending, arcp, arm_aapcs_vfpcc, arm_aapcscc, arm_apcscc, asm,
atomic, attributes, available_externally, blockaddress, builtin, byval, c,
catch, cc, ccc, cleanup, cold, coldcc, comdat, common, constant, datalayout,
declare, default, define, dereferenceable, dllexport, dllimport, eq, exact,
exactmatch, extern_weak, external, externally_initialized, false, fast, fastcc,
filter, gc, ghccc, global, hidden, inalloca, inbounds, initialexec, inlinehint,
inreg, intel_ocl_bicc, inteldialect, internal, jumptable, largest, linkonce,
linkonce_odr, localdynamic, localexec, max, min, minsize, module, monotonic,
msp430_intrcc, musttail, naked, nand, ne, nest, ninf, nnan, noalias, nobuiltin,
nocapture, noduplicate, noduplicates, noimplicitfloat, noinline, nonlazybind,
nonnull, noredzone, noreturn, nounwind, nsw, nsz, null, nuw, oeq, oge, ogt, ole,
olt, one, opaque, optnone, optsize, ord, personality, prefix, preserve_allcc,
preserve_mostcc, private, prologue, protected, ptx_device, ptx_kernel, readnone,
readonly, release, returned, returns_twice, samesize, sanitize_address,
sanitize_memory, sanitize_thread, section, seq_cst, sge, sgt, sideeffect,
signext, singlethread, sle, slt, spir_func, spir_kernel, sret, ssp, sspreq,
sspstrong, tail, target, thread_local, triple, true, type, ueq, uge, ugt, ule,
ult, umax, umin, undef, une, unnamed_addr, uno, unordered, unwind, uselistorder,
uselistorder_bb, uwtable, volatile, weak, weak_odr, webkit_jscc, x,
x86_64_sysvcc, x86_64_win64cc, x86_fastcallcc, x86_stdcallcc, x86_thiscallcc,
x86_vectorcallcc, xchg, zeroext, zeroinitializer
	},
	% types.
	%    ref: http://llvm.org/docs/LangRef.html#type-system
	keywords=[3]{
i1, i2, i3, i4, i5, i6, i7, i8, i9, i10, i11, i12, i13, i14, i15, i16, i17, i18,
i19, i20, i21, i22, i23, i24, i25, i26, i27, i28, i29, i30, i31, i32, i33, i34,
i35, i36, i37, i38, i39, i40, i41, i42, i43, i44, i45, i46, i47, i48, i49, i50,
i51, i52, i53, i54, i55, i56, i57, i58, i59, i60, i61, i62, i63, i64, i80, i512,
void, half, float, double, fp128, x86_fp80, ppc_fp128, x86_mmx, label, metadata
	},
}

% ---
% FIM DAS CUSTOMIZAÇÕES PARA A  Universidade do Estado de Santa Catarina - UDESC/CCT
% ---